% Nouveau EGA V (cote 25): archive cote 25, site page 123 (PDF page 124).
% typescript. Draft compared: none.
% AI-transcribed by Claude Opus 5.5 (claude-opus-5-5, adaptive thinking, effort high), Message Batches API, batch ega25b (msgbatch_01A7n5YuXubUPNK1zsepi8RG), 2026-11-13; checked against both supplied scan parts; see CHECK_DECISIONS.tsv.
\scanpage{25/123}
\mspage{26}

4. \uline{Systèmes linéaires de diviseurs et Modules inversibles}.

Utilisant le résultat du nº1, nous allons donner une description
complète des systèmes linéaires sur X en termes des faisceaux inversibles
sur X. Nous pouvons supposer évidemment $\check{P}\to S$ surjectif. Alors $X\times_S\check{P}\to \check{P}$ l'étant\unsure{,|.} \hand{et,} d'après les généralités du \S 20.3 la donnée d'un
diviseur  D sur $X\times_S\overset{\hand{\vee}}{P}$  revient à la
donnée d'un Module inversible $\underline{N}$ sur $X\times_S\overset{\hand{\vee}}{P}$ , et d'une section régulière
$\hand{\phi}$ de ce dernier.  L'hypothèse que D est un système linéaire de
diviseurs sur X paramétrée par $\overset{\hand{\vee}}{P}$ s'exprime alors par les deux conditions
1º) les $\hand{\phi}_t$ ($t\hand{\in}\overset{\hand{\vee}}{P}$) induits par $\hand{\phi}$ sur les fibres de $Xs_S\overset{\hand{\vee}}{P}$ au dessus de $\overset{\hand{\vee}}{P}$
sont réguliers (ce qui entraine que $\hand{\phi}$ est régulier), et 2º) les $\underline{N}_x$ ($x\hand{\in} X$)
induits par $\underline{N}$  sur les fibres des $X\times_S\overset{\hand{\vee}}{P}$ sur X sont \uline{de degré} 1 .
Or la donnée d'un $\underline{\unsure{N|E}}$ inversible sur \unsure{le} fibré projectif $X\times_S\overset{\hand{\vee}}{P}$ sur X,
satisfaisant à la condition 2º) ci-dessus, équivaut en vertu du nº1
à la donnée d'un Module inversible $\underline{L}$ sur X, $\underline{N}$ étant déterminé en fonction
de $\underline{L}$ par
\[ \underline{N}\hand{\simeq}\underline{L}\otimes_{\underline{O}_{\unsure{S|X}}}\underline{O}_{\overset{\hand{\vee}}{P}}(1) \quad , \]
et $\underline{L}$ étant d'ailleurs déterminé en termes de $\underline{N}$ par $\underline{L}\hand{\unsure{\simeq|\longrightarrow}} pr_{1\hand{*}}(\underline{N}(-1))$,
où (-1) désigne \unsure{la|le} tensorisation par $\underline{O}_{\overset{\hand{\vee}}{P}}(-1)$ sur $\underline{O}_S$ . La donnée de $\hand{\phi}$
revient à la donnée d'une section de $\underline{L}\otimes_{\underline{O}_S}\underline{O}_{\overset{\hand{\vee}}{P}}(1)$ , i.e. d'une section
\placedmargin{$\underline{O}_{X\times_S\check{P}}(1)$}
de $pr_{1\hand{*}}(\underline{L}\otimes_{\underline{O}_X}\underline{O}_{\unsure{X\times_S\check{P}}}(1))$ = $\underline{L}\otimes_{\underline{O}_S}\unsure{\illegible{1}|pr}_{\hand{*}}(\underline{O}_{\unsure{\check{P}}}(1))$, 
 or en vertu de III 2 on a $\unsure{\illegible{1}|pr}_{\hand{*}}(\underline{O}_{X\times_S\overset{\hand{\vee}}{P}}(1))\hand{\simeq}\overset{\hand{\vee}}{\underline{E}}_X$ , donc
la donnée de $\hand{\phi}$ équivaut à  la donnée d'un homomorphisme $g^{\hand{*}}(\underline{E})\hand{\to}\underline{L}$,
où ce qui revient au même, d'un morphisme u: $\underline{E}\hand{\to} g_{\hand{*}}(\underline{L})$ (NB $g:X\hand{\to} S$ est
la projection canonique). Il reste à exprimer la condition 1º) plus
haut en termes de u. Comme les constructions faites commutent avec le
changement de base, il suffit d'exprimer cette condition fibre par fibre\hand{,}
et compte tenu que les points de $\overset{\hand{\vee}}{P}$ à valeurs dans une extension k de
k(s) correspondent exactement aux droites de $E(s)\otimes_{k(s)}k$, cette

\mspage{123}
